Linear algebra is the part of mathematics that works. Whenever a problem can be made linear, it can be
solved, and a surprising number of problems that look nothing like matrices can be made linear: counting
paths that never cross, ranking the pages of the web, deciding which graphs can exist, multiplying
matrices faster than anyone thought possible, or making a learning algorithm converge in a hundred steps
instead of ten thousand.

This book takes the linear algebra of a first course to the depth that modern mathematics uses. It
proves the spectral theorem and its many relatives, the singular value decomposition, the Perron--Frobenius
theorem, the Jordan form by Young diagrams, the inequalities of Hadamard, Minkowski and Weyl, and the
theory of tensors, and it shows each one at work: in lozenge tilings and the arctic circle, in Hao Huang's
two-page proof of the sensitivity conjecture, in the codes of Delsarte, in Strassen's algorithm and in the
heavy-ball method of machine learning. Every theorem is proved, every chapter ends with exercises whose
solutions explain how to find the idea, and a Python notebook brings each chapter to life.
